\section{The actual denominator table and its uniform global export}
\label{sec:denominator-export-v050}

\status{DNE1--DNE14 reconstruct the full implications
\texttt{ConclF}$\Rightarrow$\texttt{AdmF}$\Rightarrow$
\texttt{LinkN}, including the actual local table, its coefficient-place
weights and the global prime sum.  The original common arithmetic
supplier of \texttt{ConclF} is retained.  Consequently
\texttt{GLOBAL\_DENOMINATOR} remains open.  A deliberate simplification
replaces the source's particular pointwise \texttt{Large} cutoff by a
uniform eventual cutoff; the original \texttt{LinkN} conclusion and
its constant $2000(|C|+1)^3$ are unchanged.}

The notation in this chapter always denotes the objects of one original
$D:\texttt{D7}$, $L:\texttt{SizeData}_2(D)$ and
$P:\texttt{ParameterChoice}'(D,L)$.  No free module, table on a different
host, or replacement coefficient field is introduced.  Put
\[
 x=L.\mathrm{sc}.x,\quad \epsilon=L.\mathrm{sc}.\epsilon,
 \quad z=x^\epsilon,\quad b=e^z,\quad X'=e^{z^{1/3}},
 \quad b'=xb/X',\quad w''=b/(nX').                         \tag{DNE.1}
\]
Here $n=D.R.n$, $f=2n$, $m=D.R.m$, $k=D.R.k$,
$r=k+m-1$, $t=D.R.t$, $N=k+m(t+1)$ and $n_V=r+t$.
We have $m\ge2$, $m\le f\le k$, $r\le2k$, $t\ge r+5$,
and $tr$ is odd.  Write $S=L.\mathrm{Sc}$, $\eta=L.\mathrm{etaSz}$,
$\delta=L.\mathrm{de2}$ and
\[
\begin{split}
 \mathcal D&=nt^2(N+t)z^2+kn^2(S+t+z),\\
 p_*&=mb/\mathcal D,\qquad A=kz/(m\delta),\qquad R_0=n/z,\\
 H&=b+S+\eta+t+2,\qquad \eta_3=\eta+t+2,\\
 w_B&=\max\{w''+\eta,\ \sup_{a\in D.\mathrm{host}.TT}b_a+b'/x+\eta\},\\
 w_P&=w_B+(t+1)/2,\qquad G_P=w_P+S+\eta_3,\\
 e_p&=C\left(1+H/p+\frac{\log H}{\log p}\right),\qquad
 j_p=\max\{1,\lceil\log_{p-1}A\rceil\}.
\end{split}                                                       \tag{DNE.2}
\]
The source names are respectively \texttt{Den}, \texttt{lStar},
\texttt{lA}, \texttt{R0}, \texttt{hB}, \texttt{eta3},
\texttt{wB}, \texttt{wP}, \texttt{GP}, \texttt{eP} and \texttt{jA}.
The empty supremum in $w_B$ has the source's zero convention.
The host field is $E$, with $[E:\mathbb Q]=2f$; the coefficient
field of the rows is denoted $K$, and need not be $E$.
The two generator primes form $G=\{g_0,g_1\}$, the small primes
are the divisors of $Q_S$, and $q,d$ are the two distinct ramified
auxiliary primes.  The host embeds in
$\mathbb Q(\zeta_M)$ for $M=Q_Sqd$, and $4\mid Q_S$.

\subsection{Exact local data and weights}

% BEGIN DNE1
\paragraph{DNE1: place weights and the unchanged zero-row rank.}
\label{claim:DNE1}
For a coefficient place $w$ above a rational prime $p$, the actual
weight in $\eta\mathrm{Sum}$ is
\[
 a_w=\frac{\log \mathrm Nw}{[K:\mathbb Q]}e_w
     =\frac{e_w f_w}{[K:\mathbb Q]}\log p.                 \tag{DNE.3}
\]
Thus $a_w\ge0$, and the sum over any subset of places above $p$
is at most $\log p$.  Likewise, for the host places $\lambda\mid p$,
put $n_\lambda=e_\lambda f_\lambda$.  Then
\[
 \sum_{\lambda\mid p}n_\lambda=2f,\qquad
 \sum_{\lambda\mid p}(n_\lambda+1)\le4f.                 \tag{DNE.4}
\]
These are the number-field degree formula; every residue degree
and ramification index is positive.  In the source's finite-set
presentation one may use the corresponding inequalities, which
are sufficient below.

Fix a lift $j\in\{1,\ldots,m\}$ and write $d_j=L.dZ0(j)$.
The original \texttt{adm1}, \texttt{adm2}, \texttt{adm3} say,
away from $G$, that the local row loss is bounded by
\[
 \eta_{\mathrm{loc},i}\le
    \sum_{\pi\in Z_i^{\mathrm{good}}}
         m_{i,\pi}(v_{i,\pi}+v_p(J)),\qquad
 \sum_{i\in\operatorname{locs}(j,w)}
    \sum_{\pi\in Z_i^{\mathrm{good}}}m_{i,\pi}
 \le \operatorname{rk}_0(j,w)=d_j,                       \tag{DNE.5}
\]
where $J=D.\mathrm{what}_2.\mathrm{restRowsFiniteGroupCard}>0$.
The actual equality
$\eta_J(j,w)=\sum_{i\in\operatorname{locs}(j,w)}\eta_{\mathrm{loc},i}$
is part of the $D$ record.  Consequently, if every summand at $p$
satisfies $v_{i,\pi}+v_p(J)\le U_p$, with $U_p\ge0$, then
\[
 \eta\mathrm{At}(L,j,p)\le d_j\log p\,U_p.              \tag{DNE.6}
\]
Indeed multiply (DNE.5) by the nonnegative $U_p$, use the exact rank
identity, multiply by $a_w$, and sum using (DNE.3).  This also
handles an empty place set or $d_j=0$; no division by $d_j$ occurs.
% END DNE1

% BEGIN DNE2
\paragraph{DNE2: an elementary sufficient cyclotomic budget.}
\label{claim:DNE2}
For the actual cyclotomic character $\chi_p:G_E\to\mathbb Z_p^\times$
and the odd positive integer $a=tr$, set
\[
 \nu_p=\min_{\sigma\in G_E}v_p(\chi_p(\sigma)^a-1).
\]
This minimum is the source's \texttt{nuOf}, hence its real image is
\texttt{nuL}.  For every prime $p$,
\[
 \nu_p\le v_p(M)+v_p(a),\qquad
 \sum_{p\in T}\nu_p\log p\le\log M+\log a              \tag{DNE.7}
\]
for every finite prime set $T$.

\begin{proof}
An automorphism of a cyclotomic extension whose exponent is $1$
modulo $M$ fixes $E$, because $E$ is embedded in
$\mathbb Q(\zeta_M)$.  For any prescribed finite $p$-power precision,
the Chinese remainder theorem permits every $p$-adic unit congruent
to $1$ modulo $p^{v_p(M)}$ as such an exponent.  Extend the resulting
automorphism of the finite cyclotomic compositum to the algebraic
closure over $E$.

If $p$ is odd and $p\nmid M$, choose the unit $-1$; as $a$ is odd,
$(-1)^a-1=-2$ is a unit.  If $p\mid M$, put $h=v_p(M)$ and choose
$1+p^h$.  The binomial theorem gives
\[
 v_p((1+p^h)^a-1)=h+v_p(a)
\]
for odd $p$.  One direct proof first treats an exponent prime to
$p$, for which the linear term has strictly smaller valuation than
the others, and then repeatedly raises to the $p$th power; each
step increases the valuation by one.  For $p=2$ we have $h\ge2$
and $a$ odd, so the linear term again proves the same equality.
Use precision strictly greater than this displayed valuation.
Approximation of the unit by an actual cyclotomic character value
preserves it.  This proves the first inequality and also proves
that the minimum is finite.  The second follows from
$\sum_pv_p(u)\log p=\log u$ for positive integers $u$.
\end{proof}

The source proves the sharper budget $\log Q_S+\log(tr)$ by removing
$q,d$.  The sufficient budget (DNE.7) is used here only as an
intermediate estimate: $\log M\le z/72$ follows from the original
parameter inequality $M^{36}t^{150}\le b^{1/2}$.  The final
\texttt{LinkN} bound below is exactly the original one.
% END DNE2

% BEGIN DNE3
\paragraph{DNE3: patching the actual local bounds, including the prime three.}
\label{claim:DNE3}
Suppose the original \texttt{Concl}(L,P,X,C) is supplied, where
$X$ is its actual \texttt{HubData}.  Its representations are denoted
$\rho_i$; its finite good-point sets, exceptional prime set and band
predicates are $Z_i^{\mathrm{good}}$, $\mathcal E$ and $B_{i,\lambda}$.
Put
\[
 L^X_{i,\lambda}=\max_{\pi\in Z_i^{\mathrm{good}}}
                       X.\mathrm{len0}_{i,\pi,\lambda},\qquad
 B_i^X=\max_{\pi\in Z_i^{\mathrm{good}}}
                       X.\mathrm{errorBoundHub}_{i,\pi},           \tag{DNE.8}
\]
with maximum zero for an empty set.  The same representation gives
$s_{i,\pi}=\operatorname{sPi}(\rho_i,\pi)$.
Take the one table $T_f$ provided by the \texttt{StarF} field, and set
$b_i^f=\sum_{\lambda\mid p_i}(T_f.\mathrm{len})_{i,\lambda}
             +T_f.\mathrm{BB}_i$.
Define the ordinary range
\[
 \mathcal R=\{p:p\notin\mathrm{small},\ p\ne3,q,d,
                         \ p\le p_*\ \hbox{or}\ p\in\mathcal E\}.
\]
The literal \texttt{tab} of the source is the following table,
with the branches taken in the displayed order:
\[
 \ell_{i,\lambda}=
 \begin{cases}
  0,&p_i\in\mathrm{small},\\
  b_i^f/\#\{\lambda\mid3\},&p_i=3\notin\mathrm{small},\\
  L^X_{i,\lambda},&p_i\in\mathcal R,\\
  0,&\text{otherwise},
 \end{cases}\qquad
 B_i=
 \begin{cases}
  b_i^f,&p_i\in\mathrm{small},\\
  \operatorname{midG}(L,C,p_i),&p_i=q\text{ or }d,\\
  B_i^X,&p_i\in\mathcal R,\\
  0,&\text{otherwise}.
 \end{cases}                                                   \tag{DNE.9}
\]
It retains $s$, $B_{i,\lambda}$ and $\mathcal E$, and sets
$\mathrm{gcost}(j,p)=Cx^C\eta$, $\mathrm{gs}(j,p)=p^{-2}$.
The division in (DNE.9) is legitimate: the number field $E$ has
finitely many and at least one prime above $3$.  For example $3$
is a nonunit in $\mathcal O_E$, since its norm is
$3^{[E:\mathbb Q]}$, so a maximal ideal contains it.
Thus summing its equal shares gives exactly $b_i^f$.

For every twisted local row away from $G$, (DNE.9) satisfies
\[
 v_{i,\pi}\le s_{i,\pi}+\sum_{\lambda\mid p_i}\ell_{i,\lambda}
       +V_{p_i}+v_{p_i}(h_E)+B_i+p_i^{-2}.               \tag{DNE.10}
\]
Here $V_p=D.\mathrm{what}_2.VV(p)$ and $h_E$ is the actual host
class number.  At a small prime or at $3$, this is exactly
\texttt{StarF}, by the equal-share identity.  At $q,d$, it is
\texttt{StarG}.  In $\mathcal R$, dominate each local length and
error by its maximum (DNE.8) and use \texttt{StarAD}.
In the remaining case $p>p_*$ and $p\notin\mathcal E$, so
\texttt{StarE} gives (DNE.10) with zero added length and error.
These alternatives exhaust the original prime cases.
% END DNE3

\subsection{All admissibility fields on this same table}

For precision, the numerical part of \texttt{AdmT} is recorded here.
The table has error $B_i\le B_0=Cfnt^2n_Vz^2$ away from
$\mathrm{small}\cup\{q,d\}\cup G$, and
$s_{i,\pi}+p_i^{-2}\le\nu_{p_i}+2p_i^{-2}$.
For the split and inert rows, put
\[
 F_{i,\lambda}(p)=\frac{w_P+\eta_3}{p}+S+z+
                          \eta_3\mathbf1_{B_{i,\lambda}},
 \quad R(p)=\frac H{(p-1)^{j_p+1}}+(G_P+1)(j_p+1)+z.     \tag{DNE.11}
\]
The four bounds, at the source's twisted rows away from $G$, are
\[
\begin{array}{ll}
\text{split},\ p\ge5,\ p\le p_*\text{ or }p\in\mathcal E:
 &2\sum_{\lambda\mid p}\ell_{i,\lambda}
       \le\sum_{\lambda\mid p}(2n_\lambda m+1)e_p,\\[2pt]
\text{inert},\ p<A,\ \operatorname{ord}_{\pm}(p)>4L_w:
 &\ell_{i,\lambda}\le C(n_\lambda+1)rR(p),\\[2pt]
\text{inert},\ p\ge A,\ p\le p_*\text{ or }p\in\mathcal E:
 &\ell_{i,\lambda}\le C(n_\lambda+1)r(H/p^2+F_{i,\lambda}(p))
                              +Cn_\lambda m e_p,\\[2pt]
\text{inert},\ p<A,\ \operatorname{ord}_{\pm}(p)\le4L_w:
 &\ell_{i,\lambda}\le C(n_\lambda+1)rH/p^2
       +C(n_\lambda+1)rF_{i,\lambda}(p)+Cn_\lambda m e_p.
\end{array}                                                     \tag{DNE.12}
\]
The order is in
$(\mathbb Z/2^{e+2}\mathbb Z)^\times/\{\pm1\}$.
The band implication is
\[
 B_{i,\lambda},\ \lambda\mid p\quad\Longrightarrow\quad
 \exists s\in[w'',w_B]:\ |ps-b|
   \le C\{p\eta_3+pS+b/p+s+pz\}.                       \tag{DNE.13}
\]
Above $p_*$ outside $G\cup\mathcal E$, both $\ell$ and $B_i$ vanish.
The exception bounds are
\[
 |\mathcal E|\le(|TT|+2)(S+3\eta+t+8b/p_*+9),\qquad
                    p\in\mathcal E\Longrightarrow p\le3b.       \tag{DNE.14}
\]
For $p\in\mathrm{small}\cup\{3\}$ the sum of the lengths and
error is at most $Cfm b\log x+Cx^Cz^2$; for $p=q,d$ it is at most
$\operatorname{midG}(L,C,p)\le2Cfm b$.
For $p\in G$ the direct bound is
\[
 \eta\mathrm{At}(L,j,p)\le
    d_j\log p\{Cx^C\eta+2p^{-2}+V_p+v_p(h_E)\}.        \tag{DNE.15}
\]
Finally \texttt{ParametersPolynomiallyBounded} means exactly
\[
 \begin{gathered}
  P.CLw\le C,\quad P.CYL\le C,\quad k,n,t\le x^C,
  \quad x^{-\epsilon}\le\delta,\\
  \log h_E\le Cx^C,\qquad \log J\le Cx^Cz .
 \end{gathered}                                                   \tag{DNE.16}
\]
In particular a supplied admissible table forces $C>0$, since
$P.CLw>0$.

% BEGIN DNE4
\paragraph{DNE4: the full original table is admissible.}
\label{claim:DNE4}
The table (DNE.9) satisfies all seventeen fields of
\texttt{AdmT}$(L,\mathrm{tab},C,P)$.

\begin{proof}
The first three fields are \texttt{Concl.pass},
\texttt{Concl.mult} restricted to places outside $G$, and
\texttt{Concl.rank}.  DNE3 proves the fourth field (DNE.10).
For the fifth, at the ordinary range take the maximum in
\texttt{BBLe}; outside it $B_i=0$, including $p=3$ when it is not
small.  The required upper bound is nonnegative.  For the sixth,
\texttt{KillAt} gives $s_{i,\pi}\le\nu_p+p^{-2}$, and
$\mathrm{gs}+p^{-2}=2p^{-2}$.

For the four rows (DNE.12), in $\mathcal R$ the new length is
exactly the original maximum length supplied by \texttt{RowNum}.
Outside $\mathcal R$ it is zero.  The split row excludes $3$
by $p\ge5$; the inert rows exclude it by the original
\texttt{Concl.small} clause that $3$ is not inert.  Every
right-hand side in (DNE.12) is nonnegative, since $C\ge0$,
$H\ge1$, all ranks are nonnegative, and $p\ge2$.
Thus the zero alternative also satisfies all four bounds.
The band predicate was not changed, so (DNE.13) is unchanged.

The \texttt{Small} field puts $3,q,d$ and every small prime below
$p_*$.  Hence a prime above $p_*$ outside $\mathcal E\cup G$ lies
in the zero branch of (DNE.9).  This proves large-prime vanishing;
the exceptional set itself is unchanged, giving (DNE.14).
The small-prime bound is preserved by putting its whole cost in
$B_i$, or, at $3$ outside the small set, distributing it equally
over the primes above $3$.  At $q,d$, both primes are outside the
small set and unequal to $3$, and the length is zero while the
error is exactly $\operatorname{midG}$.  The retained
\texttt{Small} bound gives its advertised upper bound.
Finally (DNE.15) is \texttt{StarH} with the chosen
$\mathrm{gcost},\mathrm{gs}$, and (DNE.16) is
\texttt{Concl.par}.  This accounts for each field on the actual
table and actual host.
\end{proof}

As a separate elementary constructor, the source's
\texttt{BoundAtTwoThree} and \texttt{StarS} imply \texttt{StarF}:
take all length entries zero and $\mathrm{BB}_i=Cfm b\log x+Cx^Cz^2$.
At $2$ use the former hypothesis; at every other small prime use
the latter; at $3$ outside the small set use the former again.
This proves the source's full \texttt{starF\_of} implication
without dropping the $2$ or $3$ case.
% END DNE4

% BEGIN DNE5
\paragraph{DNE5: a scalar majorant and the original global row sum.}
\label{claim:DNE5}
For any admissible table define $I(p)\in\{0,1\}$ by the existence
of the $s$ in (DNE.13), and put
\[
 F(p)=(w_P+\eta_3)/p+S+z+\eta_3I(p).
\]
Define
\[
 \mathcal L(p)=
 \begin{cases}
 (2fm+f)e_p,&p\text{ split},\\
 4fCrR(p),&p<A,\ \operatorname{ord}_{\pm}(p)>4L_w,\\
 4fCr(H/p^2+F(p))+2fCm e_p,&\text{otherwise},
 \end{cases}                                                     \tag{DNE.17}
\]
and the ordered scalar table
\[
 \Lambda(p)=
 \begin{cases}
 Cx^C\eta,&p\in G,\\
 Cfm b\log x+Cx^Cz^2,&p\in\mathrm{small}\text{ or }p=3,\\
 2Cfm b,&p=q\text{ or }d,\\
 0,&p>p_*,\ p\notin\mathcal E,\\
 B_0+\mathcal L(p),&\text{otherwise}.
 \end{cases}                                                     \tag{DNE.18}
\]
These are the literal source \texttt{lenB}, \texttt{lamG}; in
particular the occurrence of $C$ inside $e_p$ is retained.

For every lift $j$, assuming only summability of the original
$\mathrm{vvTerm}(D,p)$, one has
\[
\begin{split}
 \eta\mathrm{Sum}(L,j)\le d_j\biggl(&6+\log M+\log(tr)
                      +\log h_E+\log J\\
             &+\sum_{p\in T}\log p\,\Lambda(p)
                         +\sum_p\mathrm{vvTerm}(D,p)\biggr),
\end{split}                                                       \tag{DNE.19}
\]
where $T$ is exactly the image under residue characteristic of the
twisted places in $L.\mathrm{supp}$.

\begin{proof}
Sum (DNE.12) over $\lambda$, using (DNE.4).
In the split case, division by $2$ gives $(2fm+f)e_p$.
In the other cases, the bounds $\sum(n_\lambda+1)\le4f$
and $\sum n_\lambda\le2f$ give (DNE.17); the band implication
ensures $F_{i,\lambda}(p)\le F(p)$.  Together with the remaining
admissibility fields this proves
$\sum_{\lambda\mid p}\ell_{i,\lambda}+B_i\le\Lambda(p)$
outside $G$.

Use (DNE.10), the sixth field of admissibility, and DNE1, with
\[
 U_p=\nu_p+2p^{-2}+\Lambda(p)+V_p+v_p(h_E)+v_p(J).
\]
Every term here is nonnegative.  At $G$, (DNE.15) gives the same
bound directly, with harmless nonnegative extra terms.  Grouping
the original finite place sum by residue characteristic is an exact
finite-sum identity.  The cyclotomic sum is bounded by DNE2;
$\sum_{p\in T}2\log p/p^2\le4<6$ follows by comparison with the
integral of $\log u/u^2$ on $[2,\infty)$ and its first term.
Unique factorization gives
$\sum_{p\in T}v_p(h_E)\log p\le\log h_E$ and the same statement
for $J$.  Finally
\[
 \log p\,V_p=
  \log p\sum_{a=1}^{t-1}\mathrm{shaLen}_D(p,a)
       =\mathrm{vvTerm}(D,p).
\]
These are nonnegative terms.  Summability therefore permits bounding
their finite sum by the full series on this very host.  This proves
(DNE.19), retaining the exact $d_j$, $\mathrm{shaLen}$ and $\mathrm{vvTerm}$.
\end{proof}
% END DNE5

\subsection{Uniform estimates with a simpler eventual cutoff}

% BEGIN DNE6
\paragraph{DNE6: the uniform scale comparison.}
\label{claim:DNE6}
Fix $0<\epsilon<1/52$ and $C>0$.  All assertions below are uniform
over every original $D,L,P$ satisfying (DNE.16) with these fixed
$\epsilon,C$, once $x\ge x_0(\epsilon,C)$.
The original parameters give
\[
\begin{gathered}
 \eta\le e^{z^{1/4}},\quad S=4k\Delta\le4x^C e^{z^{1/4}/2},
 \quad t\le\eta,\quad 0<\delta\le1,\\
 b_a\le2b',\quad w_P+\eta_3\le9b',\quad G_P\le10b',\quad
 \eta_3\le4\eta,\quad r\eta_3\le4m w'',\\
 z\le A\le x^{C+2\epsilon},\quad
 \frac rA\le\frac{2m}{z},\quad
 p_*\le b/z,\quad \log p_*\le z,\quad A<p_*,\\
 o\le x^{\epsilon/100},\quad 2\le L_w\le C\log x,
 \quad n\le o2^e,\quad \frac km\le R_0^2,
 \quad R_0\ge\sqrt{x^{1/2}\log x}.
\end{gathered}                                                    \tag{DNE.20}
\]
Every finite polynomial with fixed coefficients in
$k,n,m,r,t,N,n_V,z,S,\eta,|TT|,\mathcal D,b/p_*$ is bounded by
\[
             Kx^a\exp(Bz^{1/4})                                \tag{DNE.21}
\]
for $K,a,B$ depending only on $C$ and that polynomial.  Such a
bound is $o(b)$, and its product with $b'$ is also $o(b)$.

\begin{proof}
The bounds on $\eta,\Delta$ are \texttt{logeta}, \texttt{logDel}.
The polynomial rank bounds and $|TT|\le2n x^{-\epsilon/8}\le2n$
give (DNE.21) by its displayed formula for $\mathcal D$ and
$b/p_*=\mathcal D/m$.
Since $z^{1/4}=x^{\epsilon/4}$ and
$z^{1/3}=x^{\epsilon/3}$, taking logarithms proves, for every
fixed $K,a,B$,
\[
 Kx^a e^{Bz^{1/4}-z}\longrightarrow0,
 \qquad Kx^{a+1}e^{Bz^{1/4}-z^{1/3}}\longrightarrow0.     \tag{DNE.22}
\]
This also proves $A<p_*$ uniformly: $A\mathcal D/(mb)$ tends to zero.

For the top weights, the \texttt{bT\_le} inequality and
\texttt{rho\_le} show $b_a\le2b'$.  If $L.\mathrm{LaG}<0$,
the correction is nonpositive; otherwise discard the nonnegative
second term of \texttt{rhoP} and bound its first term by
$x^{-2}\le1$.  The definitions then give $w_B\le3b'+\eta$.
The inequalities \texttt{teta} and \texttt{kS\_le} give
$\eta,S\le b'$, while $t\le\eta$ and $\eta\ge1$ give
$\eta_3\le4\eta$.  Substitution proves the bounds on $w_P,G_P$.
Moreover \texttt{Neta} gives
$r\eta_3\le4N\eta\le4\delta m w''\le4m w''$.
The formulas in (DNE.2), $r\le2k$ and
$\delta\in[x^{-\epsilon},1]$ give the bounds on $A,r/A$.
Since $\mathcal D\ge kn^2z\ge mz$, $p_*\le b/z$.
The final line is \texttt{o\_tie}, \texttt{Lw\_le}, the original
frame identity $n\le o2^e$, and \texttt{radius\_ge\_sqrt\_ratio}
together with \texttt{k\_ge\_mx}.  This proves every assertion.
\end{proof}

Only finitely many comparisons of type (DNE.22) are used below.
Their thresholds can therefore be maximized before any local index,
coefficient place, prime set or member of the family is chosen.
% END DNE6

% BEGIN DNE7
\paragraph{DNE7: the small-order inert primes.}
\label{claim:DNE7}
Uniformly as in DNE6, every finite set $T$ of odd integers $p\ge3$
with $\operatorname{ord}_{\pm}(p)\le4L_w$ satisfies
\[
           r\sum_{p\in T}\frac{\log p}{p^2}\le4mz^2.     \tag{DNE.23}
\]
No distribution theorem for primes is needed.

\begin{proof}
The elementary dyadic calculation proved in LFS2 says that, with
$v=\lfloor\log_2(4L_w)\rfloor$ and
$M_0=2^{e+2-v}$, these integers are $\pm1\pmod{M_0}$.
Eventually $2^e/L_w\ge2R_0\ge R_0+1$, since
$n=R_0z\le o2^e$ and $2oL_w\le z$ by (DNE.20).
Since $R_0\ge3$ eventually, $2^e\ge6L_w>4L_w$,
so $v<e$.  Thus the residue calculation applies without truncation, and
$M_0\ge2^e/L_w\ge R_0+1$.

For $R\ge3$ and an integer $M_0\ge R+1$, a decreasing-function
comparison along each of the progressions $aM_0\pm1$ gives
\[
 \sum_{p\equiv\pm1\ (M_0),\ p\ge2}\frac{\log p}{p^2}
 \le\frac{2\log R}{R^2}
             +\frac{2(\log R+2)}{RM_0}.                 \tag{DNE.24}
\]
For completeness, in each progression the first term is at most
$\log R/R^2$, because its first admissible member is at least
$M_0-1\ge R$.  The remaining terms, separated by $M_0$, are
bounded by $M_0^{-1}\int_R^\infty\log u/u^2\,du$; the latter
is $(\log R+1)/(RM_0)$, even stronger than the stated term.
Apply (DNE.24) at $R=R_0$.  Since $r\le2mR_0^2$, its first
contribution after multiplication by $r$ is at most $4m\log R_0$.
Since $n=R_0z\le oM_0L_w$, the second is at most
$4moL_w(\log R_0+2)/z$.
Now $\log R_0\le C\log x$, $o\le z^{1/100}$ and
$L_w\le C\log x$; thus
$\log R_0+oL_w(\log R_0+2)/z\le z^2$ eventually.
This proves (DNE.23) with one threshold for every finite $T$.
\end{proof}
% END DNE7

% BEGIN DNE8
\paragraph{DNE8: common error and split-place contributions below $p_*$.}
\label{claim:DNE8}
For every finite prime set $T\subset[5,p_*]$, uniformly eventually,
\[
 \sum_{p\in T}\log p\,B_0\le Cfm bz^2,
 \qquad \sum_{p\in T}\log p\,e_p\le11Cbz^2.             \tag{DNE.25}
\]

\begin{proof}
For any finite integer set below $Y\ge1$,
$\sum\log u\le Y\log Y$, $\#T\le Y$, and
$\sum\log u/u\le(1+\log Y)^2$; the last assertion follows from
$\sum_{u\le Y}1/u\le1+\log Y$.
The first bound in (DNE.25) follows from
\[
 B_0p_*z=Cfnt^2n_Vz^2\frac{mb}{\mathcal D}z
      \le Cfm bz\le Cfm bz^2,
\]
using $n_V\le N+t$.  DNE6 gives $H\le2b$, $\log H\le2z$,
$p_*\le b/z$, and $\log p_*\le z$, eventually with $z\ge1$.
Consequently
\[
 \sum_{p\in T}\log p\,e_p
 =C\left(\sum\log p+H\sum\frac{\log p}{p}
                                     +\#T\log H\right)
 \le C(b+8bz^2+2b)\le11Cbz^2.
\]
The total coefficient of $e_p$ from (DNE.17) is bounded by
$(3+2C)fm$; this keeps both the split term and the $C$ already
present in the inert term.
\end{proof}
% END DNE8

% BEGIN DNE9
\paragraph{DNE9: inert primes below $A$ with large dyadic order.}
\label{claim:DNE9}
For every finite set $T\subset[5,A)$,
\[
        \sum_{p\in T}4fCr\log p\,R(p)\le40Cfm bz^2       \tag{DNE.26}
\]
uniformly eventually.  It is permissible here to include all
integers in the indicated interval, so this bounds the needed
inert primes of large order.

\begin{proof}
By definition $(p-1)^{j_p}\ge A$, and
$j_p+1\le3+\log A$ for $p\ge5$.  Therefore
\[
 R(p)\le\frac H{A(p-1)}+(G_P+1)(3+\log A)+z.
\]
The elementary bound
$\sum_{p\le A}\log p/(p-1)\le2(1+\log A)^2$ and $r/A\le2m/z$
show that the sum of the first term, multiplied by $r$, is at
most $8mb(1+\log A)^2/z\le8mb$ eventually.
For the other terms use $G_P\le10b'$, $\log A\le(C+2)\log x$,
$A\le x^{C+2\epsilon}$ and $r\le2x^C$.
Their sum times $r$ is bounded by a fixed polynomial in $x$
times $b'$, plus a fixed polynomial in $x$.  DNE6 makes this
at most $mb$ eventually, uniformly over $T$.
Multiplication by $4fC$ gives $36Cfm b$, hence (DNE.26).
\end{proof}
% END DNE9

% BEGIN DNE10
\paragraph{DNE10: the remaining inert terms $H/p^2$.}
\label{claim:DNE10}
For the inert primes in the third and fourth rows of (DNE.12),
their combined $H/p^2$ contribution satisfies
\[
 \sum 4fCrH\frac{\log p}{p^2}\le96Cfm bz^2              \tag{DNE.27}
\]
uniformly eventually.

\begin{proof}
For $p\ge A\ge2$, monotonicity and integration of $\log u/u^2$
give the convenient bound
$\sum_{p\ge A}\log p/p^2\le4(\log A+2)/A$.
With $r/A\le2m/z$, $H\le2b$ and
$\log A+2\le z$, this contributes at most $64Cfm b$.
For $p<A$ and order at most $4L_w$, DNE7 and $H\le2b$ give
$4fC(2b)(4mz^2)=32Cfm bz^2$.  Add these bounds and use $z\ge1$.
\end{proof}
% END DNE10

% BEGIN DNE11
\paragraph{DNE11: the band contribution and the full flat term.}
\label{claim:DNE11}
For every finite prime set $T\subset[5,p_*]$, uniformly eventually,
\[
 \sum_{p\in T}4fCr\log p\,F(p)
 \le\{248C+16C(2C+1)^2\}fm bz^2.                      \tag{DNE.28}
\]

\begin{proof}
The contribution of $(w_P+\eta_3)/p$ before multiplication by
$4fC$ is at most
$9rb'(1+\log p_*)^2\le36rb'z^2\le36mbz^2$.
The last inequality follows uniformly from
$r\le2x^C$ and $b'/b=xe^{-z^{1/3}}$.
The contribution of $S+z$ is bounded by
\[
 r(S+z)p_*z
   \le2kn^2(S+t+z)\frac{mb}{\mathcal D}z
   \le2mbz\le2mbz^2.                                  \tag{DNE.29}
\]

For the band, DNE6 permits one threshold such that
$C(\eta_3+S+z)\le w''/4$.  Suppose $I(p)=1$, witnessed by
$s\ge w''$, and $p\ge2C$.  Then the right side of (DNE.13)
is at most $ps/4+Cb/p+ps/2$.  Since $Cb/p\le b/2$,
\[
 ps-b\le3ps/4+b/2,\qquad ps\le6b,
                         \qquad p\le6b/w''.             \tag{DNE.30}
\]
There are at most $6b/w''$ such integers.  Their logarithms are
at most $z$ because $p\le p_*$.  Thus
$r\eta_3\sum\log p\le(4mw'')(6b/w'')z=24mbz$.
The primes $p<2C$ contribute at most
$4mw''(2C+1)^2\le4(2C+1)^2mb$, using
$\sum_{p<2C}\log p\le(2C+1)^2$.
Together the three contributions before $4fC$ are bounded by
$\{62+4(2C+1)^2\}mbz^2$, proving (DNE.28).
\end{proof}

This argument uses the actual band implication.  It does not
replace the band set by arbitrary primes or discard its dependence
on the local row and place.
% END DNE11

% BEGIN DNE12
\paragraph{DNE12: large exceptional primes and the three special classes.}
\label{claim:DNE12}
Set $X_0=fm bz^2$.  Uniformly eventually, the part of
$\sum\log p\,\Lambda(p)$ from the large exceptional primes is
at most $X_0$, the generator-prime part at most $X_0$, the
small-prime part at most $2X_0$, and the $q,d$ part at most
$4CX_0$.  These bounds apply to arbitrary finite subsets of each class.

\begin{proof}
For the exceptions $p>p_*$, (DNE.14) and DNE6 bound
$|\mathcal E|$ by an expression of form (DNE.21), while
$p\le3b$ gives $\log p\le2z$.  Also $p>p_*>A$ and
$\log p\ge z/2$ eventually.  Here $H/p$ is bounded by
$2\mathcal D/m$, and $\log H/\log p\le4$.
Thus $e_p$ and $B_0$ satisfy (DNE.21).  In the inert branch,
$H/p^2$ does too, and
\[
 F(p)\le9b'/p_*+S+z+\eta_3
       =9x\mathcal D/(mX')+S+z+\eta_3
\]
is again of that form, even with the crude bound $I(p)\le1$.
The high-order branch below $A$ cannot occur.  Therefore the
entire exceptional sum is of form (DNE.21), hence at most
$b\le X_0$ eventually.

There are at most two generator primes, and \texttt{g\_Y},
$P.CYL\le C$ give $\log g_i\le Cz+\log2$.  Their cost is
$Cx^C\eta$ each, so their total is again of form (DNE.21)
and eventually at most $b$.
The small primes divide $Q_S$, whence
\[
 \sum_{p\in\mathrm{small}\cup\{3\}}\log p
      \le\log Q_S+\log3\le2C\log x+\log3.
\]
Multiplication by $Cfm b\log x$ is at most $X_0$ eventually,
because $(\log x)^2=o(z^2)$.  Multiplication by $Cx^Cz^2$
is at most another $b$, giving $2X_0$.
Finally $\log q,\log d\le z$ follows from
$M^{36}t^{150}\le b^{1/2}$.  At most two terms of size
$2Cfm b\log p$ give $4Cfm bz\le4CX_0$.
\end{proof}
% END DNE12

% BEGIN DNE13
\paragraph{DNE13: the complete original denominator inequality.}
\label{claim:DNE13}
For fixed $\epsilon,C$ and every sufficiently large $x$ as above,
any original admissible table gives, for every original lift $j$,
\[
 \eta\mathrm{Sum}(L,j)
 \le d_j\left[2000(C+1)^3
       \{fm bz^2+(\log x)^2L.Nc\}
                       +\sum_p\mathrm{vvTerm}(D,p)\right]        \tag{DNE.31}
\]
whenever the last series is summable.  This is exactly
\texttt{LinkN\_at}$(L,j,2000(C+1)^3)$.

\begin{proof}
Partition the ordered table (DNE.18) into its special classes,
ordinary primes below $p_*$, and exceptional primes above $p_*$.
One can also bound by the sum of the five overlapping nonnegative
majorants, so overlaps at special primes present no issue.
For ordinary primes use DNE8--DNE11.  The bound, in units of
$X_0=fm bz^2$, is
\[
 C+11C(3+2C)+40C+96C+248C+16C(2C+1)^2
                 =434C+86C^2+64C^3.                    \tag{DNE.32}
\]
DNE12 adds at most $4+4C$.
The non-prime part of (DNE.19) is at most $b\le X_0$
eventually: $\log M\le z/72$, $\log(tr)\le O_C(\log x)$,
$\log h_E\le Cx^C$ and $\log J\le Cx^Cz$.
Thus its inclusion changes the majorant to
$5+438C+86C^2+64C^3$, which is at most $2000(C+1)^3$ for $C>0$.
Substitute in (DNE.19).
The source quantity $L.Nc$ is nonnegative, by its definition
as a quotient of positive rank/scale factors and the nonnegative
thick-place count.  Adding $2000(C+1)^3(\log x)^2L.Nc$ therefore
gives exactly (DNE.31).  All multipliers use the same $d_j\ge0$.
\end{proof}

The small-prime Selmer finiteness guard is not removed: it enters
when \texttt{admitted} supplies the table under
\texttt{ShaFinOdd}.  The summability guard is used at precisely
the nonnegative infinite-sum step in DNE5.
% END DNE13

% BEGIN DNE14
\paragraph{DNE14: uniform export and the original common-family consumer.}
\label{claim:DNE14}
For every fixed $\epsilon$ and real $C_a$ one may choose
\[
                  C_{\mathrm{out}}=2000(|C_a|+1)^3             \tag{DNE.33}
\]
before the family $\mathcal F:\texttt{FamP}(\epsilon)$, so that
\texttt{AdmFC}$(\mathcal F,C_a)$ implies
\texttt{FamilyThinScalesBound}$(\mathcal F,C_{\mathrm{out}})$.
Consequently
\[
 \texttt{ConclF}(\mathcal F)\ \Longrightarrow\ 
 \texttt{AdmF}(\mathcal F)\ \Longrightarrow\ 
 \texttt{LinkN}(\mathcal F).                              \tag{DNE.34}
\]

\begin{proof}
If a family exists, its scale record at its supplied member at
$x=0$ fixes $0<\epsilon<1/52$; the family parameter $x$ and
the scale coordinate are identified at all sufficiently late
indices by its original equality.  Take the cutoff in
\texttt{AdmFC} and maximize it with $\mathcal F.x1+1$ and the
uniform cutoff of DNE6--DNE13.
At a later member satisfying \texttt{ShaFinOdd}, \texttt{admitted}
supplies an actual table with all fields of \texttt{AdmT}.
Its parameter bound forces $C_a\ge P.CLw>0$, so DNE13 applies
with $C=C_a$ and (DNE.33).  It is universal over the lift $j$,
and still has the original summability antecedent.
If $C_a\le0$, no such admissible table exists, so the same guarded
conclusion is vacuous; if there is no family the assertion is
also vacuous.  No choice of table is needed before the
\texttt{ShaFinOdd} premise.

This proves the full original \texttt{UniformExportStatement}:
$C_{\mathrm{out}}$ is independent of the family, and only the
final family threshold uses its initial cutoff.
The identity between the exported and original $\eta\mathrm{Sum}$
is definitional: both use the same twisted support, weights and
$\eta_J$.  This gives \texttt{LinkN} with precisely its original
quantifier order.  Finally, \texttt{ConclF} supplies one $C_a$
and one threshold, then an actual $X$ under \texttt{ShaFinOdd}
for each later member.  Apply DNE3--DNE4 to those very data;
this supplies \texttt{AdmFC} and proves (DNE.34).
\end{proof}

% END DNE14

\subsection{What remains upstream}

The construction of the actual common $X$ satisfying
\texttt{Concl} has not been assumed away.  In particular the
\texttt{NormChoiceFacts}, compatible Galois points and killing data,
the numerical row bounds on their actual lattices, and the original
two/three-prime and type-prime bounds must still be supplied on
the selected family.  DNE3--DNE14 now discharge the entire
table-patching, weighted aggregation and uniform export after those
inputs, and do so with the unchanged coefficient field and host.
They do not prove the closed selected-class supplier of these inputs,
or nonemptiness of the common construction class.

\paragraph{Source map.}
The exact pinned source ranges, full original contracts and file hashes
are recorded in \texttt{denominator\_\allowbreak export\_\allowbreak v50/\allowbreak provenance.\allowbreak json},
\texttt{contracts.\allowbreak json} and \texttt{source\_\allowbreak manifest.\allowbreak json}.
The main sources are \texttt{Bound\allowbreak At\allowbreak Two\allowbreak Three\allowbreak Family.\allowbreak lean},
\texttt{Admitted\allowbreak Statements\allowbreak Of\allowbreak Local\allowbreak Bounds\allowbreak Conclusion.\allowbreak lean}, its seven
table lemmas, \texttt{Twisted\allowbreak Places.\allowbreak lean},
\texttt{Mem\allowbreak Primes\allowbreak Over\allowbreak Of\allowbreak Mem.\allowbreak lean}, \texttt{NEq\allowbreak RZero\allowbreak Mul\allowbreak XPow.\allowbreak lean},
\texttt{Export\allowbreak Bound\allowbreak Of\allowbreak Admissible.\allowbreak lean},
\texttt{Theta\allowbreak Scalar\allowbreak Family\allowbreak Of\allowbreak Eigenvalue.\allowbreak lean} and
\texttt{Link\allowbreak NAlong\allowbreak Family\allowbreak Of\allowbreak Local\allowbreak Bounds.\allowbreak lean}.
The sole inherited mathematical argument used as a theorem is the
dyadic residue calculation LFS2; its range and hash are included.
All analytic comparisons used in the simplification are proved in
DNE6--DNE13.  No new external published theorem is invoked.
